% RecSys Odyssey repository-backed lecture note.
% Add figures with: \coursefigure{figures/name.svg}{Caption}
\section{Labels, features and scores}

\begin{abstract}
The model learns a target from the past; its score is an estimate, not a fact.
\end{abstract}

\subsection{Concept}
A label is the historical outcome chosen for learning: click, completion, rating or future retention. Features describe the user, item and context without using information unavailable at decision time. The model outputs a score such as a probability or expected value. Offline metrics test predictions on held-out logs; online experiments test whether the changed policy actually improves people’s outcomes.

\begin{equation*}
score(u, i, c) \approx P(label = 1 | user u, item i, context c)
\end{equation*}

\subsection{Key terms}
\begin{itemize}
\item \textbf{Label.} The outcome used as the learning target.
\item \textbf{Feature.} An input signal available when the decision is made.
\item \textbf{Score.} A model estimate used to compare candidates.
\item \textbf{Online test.} A controlled experiment that measures causal product impact.
\end{itemize}
